\chapter{Banachruimten en de fundamentele stellingen}
\label{ch:b3:banach}

De functionaalanalyse bestudeert oneindigdimensionale genormeerde
ruimten via de operatoren en functionalen die erop leven. Haar
stichtende ontdekking is dat \emph{\omterm{def:b3:complete:complete}{volledigheid}}, via de stelling
van Baire, sterke uniformiteit afdwingt: puntsgewijs begrensde
families operatoren zijn begrensd in norm (Banach--Steinhaus),
\omterm{def:b3:topology:continuity}{continue} bijecties hebben \omterm{def:b3:topology:continuity}{continue} inversen (de open afbeelding),
en grafieken sporen \omterm{def:b3:topology:continuity}{continuïteit} op (de gesloten grafiek). De
andere pijler, \emph{Hahn--Banach}, heeft helemaal geen
\omterm{def:b3:complete:complete}{volledigheid} nodig --- alleen het lemma van Zorn --- en
garandeert dat \omterm{def:b3:banach:operator}{duale ruimten} rijk genoeg zijn om elke vector te
zien. Dit hoofdstuk bewijst alle vier de stellingen en toetst ze
aan de klassieke rijruimten $\ell^p$, aan concrete berekeningen
van \omterm{def:b3:banach:operator}{duale ruimten}, en aan een werkelijk verrassende toepassing: er
bestaan \omterm{def:b3:topology:continuity}{continue} $2\pi$-periodieke functies waarvan de fourierreeks
in een punt \emph{divergeert} --- waarmee een vraag die
bachelorjaar 2 openliet ontkennend wordt beslecht.

Overal zijn $E, F$ genormeerde ruimten over $K = \R$ of $\C$; en
\emph{banachruimte}\index{banachruimte} betekent volledige
genormeerde ruimte.

\section{Begrensde operatoren; rijruimten}

\begin{definition}\label{def:b3:banach:operator}
$\mathcal L(E, F)$ noteert de ruimte van de \emph{begrensde}
(=\ \omterm{def:b3:topology:continuity}{continue}, bachelorjaar 2) lineaire afbeeldingen, met de
\emph{operatornorm}\index{operatornorm} $\vertiii T = \sup_{\norm
x \leq 1}\norm{Tx}$; die is submultiplicatief: $\vertiii{ST} \leq
\vertiii S\,\vertiii T$. De \emph{duale ruimte}\index{duale
ruimte} is $E' = \mathcal L(E, K)$.
\end{definition}

\begin{proposition}\label{prop:b3:banach:llcomplete}
Is $F$ een banachruimte, dan ook $\mathcal L(E, F)$; in het
bijzonder is $E'$ altijd een banachruimte.
\end{proposition}

\begin{proof}
Zij $(T_n)$ Cauchy voor $\vertiii\cdot$. Voor elke $x$ is
$\norm{T_nx - T_mx} \leq \vertiii{T_n - T_m}\norm x$: dus is
$(T_nx)$ Cauchy in $F$ en convergent; noem de limiet $Tx$. $T$ is
lineair (limieten van lineaire identiteiten); en overgaan tot de
limiet in $\norm{T_nx - T_mx} \leq \varepsilon\norm x$ (voor $n, m
\geq N$) geeft $\norm{T_nx - Tx} \leq \varepsilon\norm x$: dus
$T_n \to T$ in \omterm{def:b3:banach:operator}{operatornorm}, en $\vertiii T \leq \vertiii{T_N} +
\varepsilon < \infty$.
\end{proof}

\begin{definition}\label{def:b3:banach:lp}
De klassieke rijruimten\index{rijruimten $\ell^p$} (over $K$,
geïndexeerd door $\N$):
\[
\ell^p = \Bigl\{x = (x_n) : \norm x_p = \Bigl(\sum_n
\abs{x_n}^p\Bigr)^{1/p} < \infty\Bigr\}\quad (1 \leq p <
\infty),
\]
\[
\ell^\infty = \{x : \norm x_\infty = \sup\abs{x_n} < \infty\},
\]
en $c_0 = \{x : x_n \to 0\}$ met $\norm\cdot_\infty$. Dat
$\norm\cdot_p$ een norm is, volgt uit de ongelijkheid van
Minkowski, in het discrete geval precies zo bewezen als in
\cref{ch:b3:lp} (of door de eindigdimensionale ongelijkheid uit
bachelorjaar 2 te sommeren). Alle zijn banachruimten, en $c_0$ is
een gesloten deelruimte van $\ell^\infty$
(\cref{exo:b3:banach:3}).
\end{definition}

\begin{omfigure}
\begin{tikzpicture}[scale=1.5]
  \draw[->] (-1.45,0) -- (1.45,0) node[right,font=\small]
    {$x_1$};
  \draw[->] (0,-1.45) -- (0,1.45) node[above,font=\small]
    {$x_2$};
  % p = 1 diamond
  \draw[omDef, thick] (1,0) -- (0,1) -- (-1,0) -- (0,-1)
    -- cycle;
  % p = 2 circle
  \draw[omProp, thick] (0,0) circle (1);
  % p = 4 superellipse |x|^4 + |y|^4 = 1
  \draw[omThm, thick, smooth]
    plot[domain=0:90, samples=46]
    ({cos(\x)/(cos(\x)^4 + sin(\x)^4)^0.25},
     {sin(\x)/(cos(\x)^4 + sin(\x)^4)^0.25});
  \draw[omThm, thick, smooth]
    plot[domain=90:180, samples=46]
    ({cos(\x)/(cos(\x)^4 + sin(\x)^4)^0.25},
     {sin(\x)/(cos(\x)^4 + sin(\x)^4)^0.25});
  \draw[omThm, thick, smooth]
    plot[domain=180:270, samples=46]
    ({cos(\x)/(cos(\x)^4 + sin(\x)^4)^0.25},
     {sin(\x)/(cos(\x)^4 + sin(\x)^4)^0.25});
  \draw[omThm, thick, smooth]
    plot[domain=270:360, samples=46]
    ({cos(\x)/(cos(\x)^4 + sin(\x)^4)^0.25},
     {sin(\x)/(cos(\x)^4 + sin(\x)^4)^0.25});
  % p = infty square
  \draw[black!60, thick] (-1,-1) rectangle (1,1);
  \node[omDef, font=\small] at (0.40,0.40) {$p{=}1$};
  \node[omProp, font=\small] at (-0.52,0.75) {$p{=}2$};
  \node[omThm, font=\small] at (0.90,-0.52) {$p{=}4$};
  \node[black!60, font=\small] at (-1.62,1.15)
    {$p{=}\infty$};
  \draw[black!60, ->] (-1.32,1.11) -- (-1.02,1.0);
\end{tikzpicture}

{\small Eenheidsballen van de $p$-normen in het vlak, genest
naarmate $p$ groeit van $1$ (ruit) via $2$ (schijf) en $4$
(superellips) naar $\infty$ (vierkant). De convexiteit van elke
bal \emph{is} de ongelijkheid van Minkowski; en in de hoeken bij
$p = 1$ en $p = \infty$ ontaarden strikte convexiteit, de
eenduidigheid van beste benaderingen en de gelijkheidsgevallen
van \cref{exo:b3:lp:12} tegelijk.}
\end{omfigure}

\begin{proposition}[Neumannreeks]\label{prop:b3:banach:neumann}
Zij $E$ een banachruimte en $T \in \mathcal L(E) = \mathcal
L(E,E)$ met $\vertiii T < 1$. Dan is $I - T$ inverteerbaar in
$\mathcal L(E)$, met $(I - T)^{-1} = \sum_{n\geq0}T^n$
(convergent in \omterm{def:b3:banach:operator}{operatornorm}). Bijgevolg is de verzameling
inverteerbare operatoren open, en is het inverteren erop
\omterm{def:b3:topology:continuity}{continu}.\index{Neumannreeks}
\end{proposition}

\begin{proof}
De reeks convergeert absoluut ($\vertiii{T^n} \leq \vertiii T^n$,
meetkundig) in de banachruimte $\mathcal L(E)$
(\cref{prop:b3:banach:llcomplete};
\cref{exo:b3:complete:1}(b)). Telescoperen geeft $(I -
T)\sum_{n \leq N}T^n = I - T^{N+1} \to I$, en net zo aan de andere
kant. Voor de openheid: is $S$ inverteerbaar en $\vertiii{H} <
1/\vertiii{S^{-1}}$, dan is $S + H = S(I + S^{-1}H)$ met
$\vertiii{S^{-1}H} < 1$: inverteerbaar. \omterm{def:b3:topology:continuity}{Continuïteit} van het
inverteren: de reeksuitdrukking geeft lokaal
$\vertiii{(S+H)^{-1} - S^{-1}} = O(\vertiii H)$.
\end{proof}

\begin{example}[Een vergelijking van Volterra, met Neumann
opgelost]
\label{ex:b3:banach:volterraexample}
Beschouw op $E = \mathcal C(\intcc01)$ de integraalvergelijking
\[
u(x) = 1 + \lambda\int_0^xu(t)\,\dd t,
\qquad\text{dat wil zeggen}\qquad u = \mathbf 1 + \lambda Tu,
\quad (Tu)(x) = \int_0^xu .
\]
Hier is $\vertiii T \leq 1$, dus voor $\abs\lambda < 1$ is de
\omterm{prop:b3:banach:neumann}{Neumannreeks} rechtstreeks van toepassing: $u = (I - \lambda
T)^{-1}\mathbf 1 = \sum_n\lambda^nT^n\mathbf 1$. Berekenen geeft
$T^n\mathbf 1 = \frac{x^n}{n!}$, dus
\[
u(x) = \sum_{n\geq0}\frac{(\lambda x)^n}{n!} =
\eu^{\lambda x} ,
\]
zoals differentiëren bevestigt. Sterker nog: $\vertiii{T^n} \leq
\frac1{n!}$ (de herhaalde kern krimpt faculteitsgewijs), zodat
$\sum\lambda^nT^n$ voor \emph{elke} $\lambda$ convergeert --- de
operator $I - \lambda T$ is inverteerbaar voor alle $\lambda \in
\C$, ook al is $\vertiii{\lambda T} \geq 1$ vanaf zeker moment:
wat telt is het spectrale afnemen van de machten, niet de eerste
norm. Volterra-operatoren dragen dat faculteitsgewijze afnemen in
zich (\cref{exo:b3:complete:4}(a) buitte precies dat uit), en
daarom kennen beginwaardeproblemen nooit de
resonantieverschijnselen van randwaardeproblemen
(\cref{ch:b3:spectral}).
\end{example}

\section{Hahn--Banach}

\begin{theorem}[Hahn--Banach, analytische vorm]
\label{thm:b3:banach:hahnbanach}
Zij $E$ een \emph{reële} vectorruimte, $p \colon E \to \R$
\emph{sublineair} ($p(x + y) \leq p(x) + p(y)$ en $p(tx) = tp(x)$
voor $t \geq 0$), $F \subseteq E$ een deelruimte en $f \colon F
\to \R$ lineair met $f \leq p$ op $F$. Dan zet $f$ zich voort tot
een lineaire $\tilde f \colon E \to \R$ met $\tilde f \leq p$ op
$E$.\index{stelling van Hahn--Banach}
\end{theorem}

\begin{proof}
\emph{Voortzetting in één stap.} Zij $x_0 \notin F$; we zetten $f$
voort tot $F \oplus \R x_0$ door $\alpha = \tilde f(x_0)$ goed te
kiezen: we hebben nodig, voor alle $y \in F$ en $t > 0$,
\[
f(y) + t\alpha \leq p(y + tx_0)
\quad\text{en}\quad
f(y) - t\alpha \leq p(y - tx_0),
\]
wat na deling door $t$ (sublineariteit) neerkomt op
\[
\sup_{v \in F}\ \bigl[f(v) - p(v - x_0)\bigr]
\;\leq\; \alpha \;\leq\;
\inf_{u \in F}\ \bigl[p(u + x_0) - f(u)\bigr].
\]
Zo'n $\alpha$ bestaat precies wanneer elk linkerlid $\leq$ elk
rechterlid is: en inderdaad is $f(v) + f(u) = f(u + v) \leq p(u +
v) \leq p(u + x_0) + p(v - x_0)$, dat wil zeggen $f(v) - p(v -
x_0) \leq p(u + x_0) - f(u)$.

\emph{Zorn.} Orden de voortzettingen van $f$ die door $p$
gedomineerd worden (paren: deelruimte, functionaal) naar
voortzetting; een keten heeft de vereniging als bovengrens; en een
maximaal element moet op heel $E$ gedefinieerd zijn, want anders
spreekt de voortzetting in één stap de maximaliteit tegen.
\end{proof}

\begin{corollary}\label{cor:b3:banach:hbnormed}
Zij $E$ een genormeerde ruimte ($K = \R$ of $\C$).
\begin{enumerate}
  \item Elke $f \in F'$ (met $F$ een deelruimte) zet zich voort
        tot een $\tilde f \in E'$ met $\norm{\tilde f}_{E'} =
        \norm f_{F'}$.
  \item Voor elke $x \neq 0$ is er een $f \in E'$ met $\norm f =
        1$ en $f(x) = \norm x$. In het bijzonder scheidt $E'$ de
        punten van $E$, en is $\norm x = \sup_{\norm f \leq
        1}\abs{f(x)}$.
  \item Voor een gesloten deelruimte $F$ en $x \notin F$ is er een
        $f \in E'$ die op $F$ verdwijnt, met $f(x) = d(x, F)$ en
        $\norm f \leq 1$.
\end{enumerate}
\end{corollary}

\begin{proof}
(1) Reëel geval: pas \cref{thm:b3:banach:hahnbanach} toe met
$p(x) = \norm f_{F'}\,\norm x$ (sublineair); de voortzetting
voldoet aan $\pm\tilde f(x) = \tilde f(\pm x) \leq p(x)$, dus
$\norm{\tilde f} \leq \norm f$, en $\geq$ is de beperking.
Complex geval: zij $u = \operatorname{Re}f$, een reële
functionaal met $\abs u \leq \norm f\norm\cdot$; merk op dat $f(x)
= u(x) - \iu\,u(\iu x)$ (controleer op reëel en imaginair deel:
$\operatorname{Im}f(x) = -\operatorname{Re}f(\iu x)$). Zet $u$
reëel-lineair voort met dezelfde grens, en stel $\tilde f(x) =
\tilde u(x) - \iu\tilde u(\iu x)$: die is $\C$-lineair
(rechtstreeks na te gaan op vermenigvuldiging met $\iu$) en zet
$f$ voort; norm: schrijf voor gegeven $x$ de waarde $\tilde f(x) =
r\eu^{\iu\theta}$, dan is $\abs{\tilde f(x)} = \tilde
f(\eu^{-\iu\theta}x) = \tilde u(\eu^{-\iu\theta}x) \leq \norm
f\,\norm x$.

(2) Definieer op $F = Kx$: $f(tx) = t\norm x$, van norm $1$ op
$F$; zet voort met (1). De dualiteitsformule: $\leq$ is duidelijk,
en $\geq$ volgt uit deze $f$.

(3) Definieer op $F \oplus Kx$: $f(y + tx) = t\,d(x, F)$; dan is
voor $t \neq 0$ $\norm{y + tx} = \abs t\,\norm{x + y/t} \geq \abs
t\,d(x, F) = \abs{f(y + tx)}$, dus $\norm f \leq 1$ op die
deelruimte; zet voort met (1).
\end{proof}

\begin{remark}\label{rem:b3:banach:bidual}
Volgens (2) is de kanonieke afbeelding $J \colon E \to E''$,
$J(x)(f) = f(x)$, een isometrie (\cref{exo:b3:banach:10}): elke
genormeerde ruimte zit in haar biduale ruimte. Ruimten waarvoor
$J$ surjectief is, heten
\emph{reflexief}\index{reflexieve ruimte}; de weekendopgave toont
aan dat $\ell^p$ ($1 < p < \infty$) \omterm{rem:b3:banach:bidual}{reflexief} is en $\ell^1$
niet.
\end{remark}

\section{De trilogie van Baire}

\begin{theorem}[Banach--Steinhaus, uniforme begrensdheid]
\label{thm:b3:banach:banachsteinhaus}
Zij $E$ een \emph{banachruimte}, $F$ genormeerd, en
$(T_i)_{i\in I} \subseteq \mathcal L(E, F)$ een familie met
$\sup_i \norm{T_ix} < \infty$ voor elke $x \in E$. Dan is $\sup_i
\vertiii{T_i} < \infty$.\index{stelling van Banach--Steinhaus}
\end{theorem}

\begin{proof}
De verzamelingen $F_n = \{x : \sup_i\norm{T_ix} \leq n\}$ zijn
gesloten (doorsneden van originelen van gesloten ballen) en
overdekken $E$. Baire (\cref{thm:b3:complete:baire}) geeft een
$n_0$ en een bal $B(x_0, r) \subseteq F_{n_0}$. Voor $\norm z <
r$ is $\norm{T_iz} \leq \norm{T_i(x_0 + z)} + \norm{T_ix_0} \leq
2n_0$, dus $\vertiii{T_i} \leq 2n_0/r$ voor elke $i$.
\end{proof}

\begin{corollary}\label{cor:b3:banach:bslimits}
Is $E$ een banachruimte en convergeren $T_n \in \mathcal L(E,F)$
\emph{puntsgewijs} ($T_nx \to Tx$ voor elke $x$), dan is $\sup_n
\vertiii{T_n} < \infty$, $T \in \mathcal L(E, F)$ en $\vertiii T
\leq \liminf \vertiii{T_n}$.
\end{corollary}

\begin{proof}
Convergente rijen zijn begrensd: puntsgewijze begrensdheid;
Banach--Steinhaus begrenst de normen door zekere $M$; en dan is
$\norm{Tx} = \lim\norm{T_nx} \leq M\norm x$ ($T$ is lineair als
puntsgewijze limiet), en de scherpere grens volgt door in
$\norm{T_nx} \leq \vertiii{T_n}\norm x$ tot de $\liminf$ over te
gaan.
\end{proof}

\begin{theorem}[Divergente fourierreeksen]
\label{thm:b3:banach:fourierdiverge}
Er bestaan \omterm{def:b3:topology:continuity}{continue} $2\pi$-periodieke functies $f$ waarvan de
fourierreeks in $0$ divergeert: $\sup_N\abs{S_N(f)(0)} = \infty$.
Zulke $f$ vormen zelfs een \omterm{def:b3:topology:interior}{dichte} deelverzameling van $\mathcal
C(S^1)$.\index{fourierreeks, divergentie}
\end{theorem}

\begin{proof}
Werk in $E = (\mathcal C(S^1), \norm\cdot_\infty)$, een
banachruimte, met de functionalen $\Lambda_N(f) = S_N(f)(0) =
\frac1{2\pi}\int_{-\pi}^{\pi}f(t)\,D_N(t)\,\dd t$ (de
dirichletkern, bachelorjaar 2). Elke $\Lambda_N$ is \omterm{def:b3:topology:continuity}{continu} met
\[
\norm{\Lambda_N} = \frac1{2\pi}\int_{-\pi}^\pi\abs{D_N(t)}\,\dd
t \;=\; L_N .
\]
($\leq$ is duidelijk; $\geq$: neem een \omterm{def:b3:topology:continuity}{continue} $f$ met $\norm
f_\infty \leq 1$ die $\operatorname{sign}D_N$ benadert --- het
teken heeft eindig veel sprongen, en elke sprong afvlakken over
een interval van lengte $\varepsilon$ verandert de integraal met
$O(N\varepsilon)$.) De \emph{lebesgueconstanten} $L_N$ gaan naar
oneindig:
\[
L_N = \frac1{2\pi}\int_{-\pi}^{\pi}
\frac{\abs{\sin\bigl((N{+}\tfrac12)t\bigr)}}{\abs{\sin(t/2)}}
\,\dd t
\geq \frac{2}{\pi}\int_0^{\pi}
\frac{\abs{\sin\bigl((N{+}\tfrac12)t\bigr)}}{t}\,\dd t
= \frac{2}{\pi}\int_0^{(N+\frac12)\pi}\frac{\abs{\sin
u}}{u}\,\dd u
\]
(met $\abs{\sin(t/2)} \leq t/2$ op $[0, \pi]$, en daarna de
substitutie $u = (N + \tfrac12)t$). Snijden in bogen:
\[
\int_{(k-1)\pi}^{k\pi}\frac{\abs{\sin u}}u\,\dd u
\geq \frac1{k\pi}\int_{(k-1)\pi}^{k\pi}\abs{\sin u}\,\dd u =
\frac{2}{k\pi},
\qquad\text{dus}\qquad
L_N \geq \frac{4}{\pi^2}\sum_{k=1}^N\frac1k
\xrightarrow[N\to\infty]{} \infty .
\]
Had elke \omterm{def:b3:topology:continuity}{continue} $f$ dat $\sup_N \abs{\Lambda_N(f)} < \infty$,
dan dwong Banach--Steinhaus af dat $\sup_N\norm{\Lambda_N} <
\infty$: tegenspraak. Dus heeft een zekere $f$ --- zelfs een
\omterm{rem:b3:complete:meagre}{niet-magere}, \omterm{def:b3:topology:interior}{dichte verzameling} $f$ (het complement van
$\bigcup_M\{f: \sup_N\abs{\Lambda_Nf}\leq M\}$, een aftelbare
vereniging gesloten verzamelingen die wegens het bovenstaande,
toegepast in elke bal, geen inwendige hebben en dus \omterm{rem:b3:complete:meagre}{mager} is) ---
dat $\sup_N\abs{S_Nf(0)} = \infty$.
\end{proof}

\begin{theorem}[Open afbeelding]\label{thm:b3:banach:openmapping}
Zij $E, F$ banachruimten en $T \in \mathcal L(E, F)$
\emph{surjectief}. Dan is $T$ open: $T(B_E(0,1)) \supseteq B_F(0,
c)$ voor zekere $c > 0$. Bijgevolg heeft een bijectieve begrensde
operator tussen banachruimten een begrensde
inverse.\index{open-afbeeldingsstelling}
\end{theorem}

\begin{proof}
Schrijf $B = B_E(0,1)$. De surjectiviteit geeft $F = \bigcup_n
\overline{T(nB)} = \bigcup_n n\,\overline{T(B)}$; Baire
(\cref{thm:b3:complete:baire}) geeft $\overline{T(B)}$ een
inwendige: een zekere $B_F(y_0, 4c) \subseteq \overline{T(B)}$.
Hercentreer in $0$: voor $\norm y < 4c$ zijn zowel $y_0 + y$ als
$y_0$ limieten van beelden $Tu_k$ respectievelijk $Tv_k$ met $u_k,
v_k \in B$, dus $y = \lim T(u_k - v_k)$ met $u_k - v_k \in 2B$:
dus $B_F(0, 4c) \subseteq \overline{T(2B)}$, oftewel $B_F(0, 2c)
\subseteq \overline{T(B)}$.

\emph{De \omterm{def:b3:topology:interior}{afsluiting} weghalen} (hier komt de \omterm{def:b3:complete:complete}{volledigheid} van $E$
binnen): zij $\norm y < c$. Kies $x_1 \in \frac12 B$ met $\norm{y
- Tx_1} < c/2$ ($B_F(0,2c) \subseteq \overline{T(B)}$, met
$\frac12$ geschaald); inductief $x_k \in 2^{-k}B$ met $\norm{y -
T(x_1 + \dots + x_k)} < c\,2^{-k}$. De reeks $\sum x_k$
convergeert absoluut in de banachruimte $E$, naar een $x \in B$
(norm $< \sum 2^{-k} = 1$), en $Tx = y$ wegens de \omterm{def:b3:topology:continuity}{continuïteit}:
dus $B_F(0, c) \subseteq T(B)$. De openheid van $T$ op
willekeurige \omterm{def:b3:topology:topology}{open verzamelingen} volgt door translatie en
schaling; en voor het gevolg: dat $T$ open is, betekent dat
$T^{-1}$ \omterm{def:b3:topology:continuity}{continu} is.
\end{proof}

\begin{corollary}[Gelijkwaardige normen]\label{cor:b3:banach:equivnorms}
Is een vectorruimte \omterm{def:b3:complete:complete}{volledig} voor twee vergelijkbare normen
($\norm\cdot_a \leq C\norm\cdot_b$), dan zijn de normen
gelijkwaardig.
\end{corollary}

\begin{proof}
De identiteit $(E, \norm\cdot_b) \to (E, \norm\cdot_a)$ is
begrensd en bijectief tussen banachruimten: haar inverse is
begrensd.
\end{proof}

\begin{theorem}[Gesloten grafiek]\label{thm:b3:banach:closedgraph}
Zij $E, F$ banachruimten en $T \colon E \to F$ lineair. Is de
grafiek $\Gamma = \{(x, Tx)\}$ gesloten in $E \times F$ (dat wil
zeggen: uit $x_n \to x$ en $Tx_n \to y$ volgt $y = Tx$), dan is
$T$ begrensd.\index{gesloten-grafiekstelling}
\end{theorem}

\begin{proof}
$E \times F$ met $\norm{(x,y)} = \norm x + \norm y$ is een
banachruimte; $\Gamma$, een gesloten deelruimte, is er ook een. De
projectie $\pi_E\colon \Gamma \to E$ is begrensd en bijectief, dus
is haar inverse $x \mapsto (x, Tx)$ begrensd
(\cref{thm:b3:banach:openmapping}): $\norm{Tx} \leq \norm{(x,
Tx)} \leq C\norm x$.
\end{proof}

\begin{method}\label{met:b3:banach:trilogy}
Wanneer welke stelling. \emph{Hahn--Banach}: om een functionaal
met voorgeschreven gedrag te maken (een vector normeren, op een
deelruimte verdwijnen, vanaf een deelruimte voortzetten) --- geen
\omterm{def:b3:complete:complete}{volledigheid} nodig. \emph{Banach--Steinhaus}: om puntsgewijze
informatie in uniforme grenzen om te zetten --- meestal om te
tonen dat een limietbewerking \omterm{def:b3:topology:continuity}{continu} is, of (contrapositief) om
\emph{divergentie} voor een element te bewijzen, zoals bij de
fourierreeksen. \emph{Open afbeelding / gesloten grafiek}: om
\omterm{def:b3:topology:continuity}{continuïteit} gratis te krijgen uit \omterm{def:b3:galois:algebraic}{algebraïsche} bijectiviteit of
uit een afsluitingseigenschap van de grafiek --- typisch gebruik:
twee volledige normen vergelijken, of automatische \omterm{def:b3:topology:continuity}{continuïteit}
bewijzen. Alle drie de bairestellingen vergen \omterm{def:b3:complete:complete}{volledigheid}
\emph{van het vertrekpunt}; anders zijn er tegenvoorbeelden
(\cref{exo:b3:banach:7}).
\end{method}

\section{Duale ruimten, concreet}

\begin{theorem}\label{thm:b3:banach:duals}
Isometrisch geldt: $(c_0)' \cong \ell^1$ en $(\ell^1)' \cong
\ell^\infty$, via de paring $\langle x, y\rangle = \sum_n
x_ny_n$.\index{duale ruimten van rijruimten}
\end{theorem}

\begin{proof}
We bewijzen $(c_0)' \cong \ell^1$; de tweede identificatie is
\cref{exo:b3:banach:5}. Ken aan $y \in \ell^1$ de functionaal
$\Lambda_y(x) = \sum x_ny_n$ toe (voor $x \in c_0$): absoluut
convergent, met $\abs{\Lambda_y(x)} \leq \norm x_\infty\norm
y_1$, dus $\norm{\Lambda_y} \leq \norm y_1$. Zij omgekeerd
$\Lambda \in (c_0)'$; stel $y_n = \Lambda(e_n)$ (met $e_n$ de
eenheidsrijen). Toets voor elke $N$ op $x^{(N)} = \sum_{n \leq N}
\operatorname{sign}(\overline{y_n})\,e_n \in c_0$ (norm $\leq 1$;
gebruik in het complexe geval de factoren $\bar y_n/\abs{y_n}$ van
modulus $1$): $\Lambda(x^{(N)}) = \sum_{n\leq N}\abs{y_n} \leq
\norm\Lambda$. Dus $y \in \ell^1$ met $\norm y_1 \leq
\norm\Lambda$. Ten slotte is $\Lambda = \Lambda_y$: beide vallen
samen op de $e_n$ en dus op de eindige rijen, die \omterm{def:b3:topology:interior}{dicht} liggen in
$c_0$ (afknotten: $\norm{x - \sum_{n \leq N}x_ne_n}_\infty =
\sup_{n > N}\abs{x_n} \to 0$, juist omdat $x_n \to 0$); en
\omterm{def:b3:topology:continuity}{continue} functionalen die op een \omterm{def:b3:topology:interior}{dichte verzameling} samenvallen,
zijn gelijk. De correspondentie is lineair, bijectief en
isometrisch ($\norm{\Lambda_y} = \norm y_1$ uit de twee
ongelijkheden).
\end{proof}

\section{Oefeningen}

\begin{exercise}[$\star$]\label{exo:b3:banach:1}
Bereken de \omterm{def:b3:banach:operator}{operatornormen}:
(a) van de verschuivingen $S(x_1, x_2, \dots) = (0, x_1, x_2,
\dots)$ en $S^*(x_1, x_2, \dots) = (x_2, x_3, \dots)$ op
$\ell^2$;
(b) van de vermenigvuldigingsoperator $M_a x = (a_nx_n)$ op
$\ell^2$, voor $a \in \ell^\infty$;
(c) van de functionaal $\Lambda(f) = \int_0^{1/2}f -
\int_{1/2}^1f$ op $\mathcal C(\intcc01)$ --- toon aan dat
$\norm\Lambda = 1$ en dat de norm \emph{niet bereikt} wordt.
\end{exercise}

\begin{exercise}[$\star$]\label{exo:b3:banach:2}
Zij $E$ een banachruimte, $T \in \mathcal L(E)$ inverteerbaar, en
$S$ met $\vertiii{S - T} < 1/\vertiii{T^{-1}}$. Toon aan dat $S$
inverteerbaar is en schat $\vertiii{S^{-1} - T^{-1}}$.
Toepassing: is een lineair stelsel $Tx = b$ \omterm{def:b3:groups:derived}{oplosbaar} met $T$
inverteerbaar, dan blijft het bij een voldoende kleine verstoring
van $T$ eenduidig \omterm{def:b3:groups:derived}{oplosbaar}, met een kwantitatieve grens op de
verandering van de oplossing.
\end{exercise}

\begin{exercise}[$\star\star$]\label{exo:b3:banach:3}
(a) Bewijs dat $\ell^1$, $\ell^\infty$ en $c_0$ banachruimten
zijn, en dat $c_0$ de \omterm{def:b3:topology:interior}{afsluiting} in $\ell^\infty$ is van de
ruimte van de eindige rijen.
(b) Toon aan dat $\ell^p \subseteq \ell^q$ met $\norm\cdot_q \leq
\norm\cdot_p$ voor $1 \leq p \leq q \leq \infty$, en dat de
inclusie strikt is.
\end{exercise}

\begin{exercise}[$\star\star$]\label{exo:b3:banach:4}
Zij $F \subseteq E$ een gesloten deelruimte en $x \notin F$.
Bewijs met \cref{cor:b3:banach:hbnormed} de dualiteitsformule
\[
d(x, F) = \max\bigl\{\abs{f(x)} : f \in E',\ \norm f \leq 1,\
f\restriction_F = 0\bigr\}
\]
(let op: een \emph{maximum}). Leid af dat $F = \bigcap\{\ker f : f
\in E',\ f\restriction_F = 0\}$: gesloten deelruimten zijn precies
de doorsneden van kernen van functionalen.
\end{exercise}

\begin{exercise}[$\star\star$]\label{exo:b3:banach:5}
Bewijs isometrisch dat $(\ell^1)' \cong \ell^\infty$, volgens het
schema van \cref{thm:b3:banach:duals} (de eindige rijen liggen
\omterm{def:b3:topology:interior}{dicht} in $\ell^1$). Waar breekt het argument voor
$(\ell^\infty)'$?
\end{exercise}

\begin{exercise}[$\star\star$]\label{exo:b3:banach:6}
Zij $E, F, G$ genormeerd met $E$ een banachruimte, en $B \colon E
\times F \to G$ bilineair en in elke veranderlijke afzonderlijk
\omterm{def:b3:topology:continuity}{continu}. Toon aan dat $B$ (gezamenlijk) \omterm{def:b3:topology:continuity}{continu} is: $\norm{B(x,y)}
\leq C\norm x\norm y$. \emph{(Pas Banach--Steinhaus toe op de
familie $(B(\cdot, y))_{\norm y \leq 1}$.)}
\end{exercise}

\begin{exercise}[$\star\star$]\label{exo:b3:banach:7}
(a) Vergelijk op $E = \mathcal C(\intcc01)$ de normen
$\norm\cdot_\infty$ en $\norm\cdot_1$: de identiteit $(E,
\norm\cdot_\infty) \to (E, \norm\cdot_1)$ is begrensd en
bijectief, maar haar inverse is onbegrensd. Welke hypothese van
\cref{cor:b3:banach:equivnorms} faalt?
(b) Geef een discontinue lineaire afbeelding van een \omterm{def:b3:topology:interior}{dichte}
deelruimte van $\ell^2$ naar $K$ (bijvoorbeeld op de eindige
rijen), en leg uit waarom dat de gesloten-grafiekstelling niet
tegenspreekt.
\end{exercise}

\begin{exercise}[$\star\star\star$]\label{exo:b3:banach:8}
(Hellinger--Toeplitz) Zij $T \colon \ell^2 \to \ell^2$ lineair
(overal gedefinieerd) en \emph{symmetrisch}: $\langle Tx, y\rangle
= \langle x, Ty\rangle$ voor alle $x, y$, met $\langle x, y
\rangle = \sum x_n\bar y_n$. Toon aan dat $T$ begrensd is.
\emph{(Gesloten grafiek: als $x_k \to x$ en $Tx_k \to z$, toets
dan tegen een willekeurige $y$.)} Moraal: onbegrensde
symmetrische operatoren --- de hamiltonianen van de
kwantummechanica --- kunnen nooit op de hele ruimte gedefinieerd
worden.
\end{exercise}

\begin{exercise}[$\star\star\star$]\label{exo:b3:banach:9}
(Stelling van Pólya over kwadratuur) Zij voor elke $n$ de regel
$\Lambda_n(f) = \sum_{i=0}^{n} w_{i,n}f(x_{i,n})$ een
kwadratuurregel op $\mathcal C(\intcc01)$ (met $x_{i,n} \in
\intcc01$ en $w_{i,n} \in \R$). Toon aan dat $\Lambda_n(f) \to
\int_0^1f$ voor \emph{elke} \omterm{def:b3:topology:continuity}{continue} $f$ dan en slechts dan als:
(i) $\Lambda_n(P) \to \int_0^1P$ voor elke veelterm $P$, en (ii)
$\sup_n\sum_i\abs{w_{i,n}} < \infty$. \emph{(Bereken
$\norm{\Lambda_n}$; gebruik Banach--Steinhaus en Weierstrass.)}
Ga na dat regels met \emph{positieve} gewichten die exact zijn op
de constanten, automatisch aan (ii) voldoen.
\end{exercise}

\begin{exercise}[$\star\star$]\label{exo:b3:banach:10}
Toon aan dat $J \colon E \to E''$, $J(x)(f) = f(x)$, een lineaire
isometrie is (gebruik \cref{cor:b3:banach:hbnormed}(2)), en dat ze
surjectief is wanneer $\dim E < \infty$. Toon ook aan dat $E$
separabel is zodra $E'$ dat is. \emph{(Kies $x_n$ die een \omterm{def:b3:topology:interior}{dichte}
rij van $E'$ bijna normeren, en toon aan dat hun gesloten
opspansel $E$ is, via \cref{cor:b3:banach:hbnormed}(3).)}
\end{exercise}

\begin{exercise}[$\star\star$]\label{exo:b3:banach:11}
(Quotiëntruimten) Zij $E$ een banachruimte en $F \subseteq E$ een
\emph{gesloten} deelruimte. Definieer op $E/F$
\[
\norm{\bar x} \;=\; d(x, F) = \inf_{y\in F}\norm{x - y} .
\]
(a) Toon aan dat dit een goed gedefinieerde \emph{norm} op $E/F$
is (waar komt de geslotenheid van $F$ binnen?), en dat de
projectie $\pi \colon E \to E/F$ voldoet aan $\vertiii\pi \leq 1$
en de open eenheidsbal op de open eenheidsbal afbeeldt.
(b) Toon aan dat $E/F$ \omterm{def:b3:complete:complete}{volledig} is. \emph{(Gebruik het
reekscriterium van \cref{exo:b3:complete:1}(b): gegeven klassen
$\bar x_k$ met $\sum\norm{\bar x_k} < \infty$, til elke op tot
$x_k \in E$ met $\norm{x_k} \leq \norm{\bar x_k} + 2^{-k}$ en
sommeer in $E$.)}
(c) Bereken: toon voor $E = c$ (de convergente rijen) en $F =
c_0$ aan dat $c/c_0 \cong K$ isometrisch, via $\bar x \mapsto
\lim_nx_n$.
\end{exercise}

\begin{exercise}[$\star\star$]\label{exo:b3:banach:12}
(Begrensde projecties en gecomplementeerde deelruimten) Zij $E$
een banachruimte en $P \colon E \to E$ lineair met $P^2 = P$ (een
\omterm{def:b3:galois:algebraic}{algebraïsche} projectie), $V = \operatorname{im}P$ en $W = \ker
P$.
(a) Stel dat $P$ begrensd is. Toon aan dat $V$ en $W$ gesloten
zijn en dat $E = V \oplus W$ met de ontbinding $x = Px + (x -
Px)$.
(b) Stel omgekeerd dat $E = V \oplus W$ met $V$ en $W$
\emph{beide} gesloten, en zij $P$ de projectie op $V$ langs $W$.
Toon aan dat $P$ begrensd is. \emph{(Gesloten grafiek: als $x_n
\to x$ en $Px_n \to z$, dan is $z \in V$ en $x_n - Px_n \to x - z
\in W$, en de eenduidigheid van de ontbinding geeft $z = Px$.)}
(c) Leid de \emph{gelijkwaardigheid} af: een deelruimte $V$ laat
een begrensde projectie toe dan en slechts dan als ze gesloten is
en een gesloten \omterm{def:b3:galois:algebraic}{algebraïsch} complement heeft --- en merk op
(zonder bewijs) dat er gesloten deelruimten zonder die eigenschap
bestaan ($c_0$ binnen $\ell^\infty$ is het klassieke voorbeeld):
hilbertruimten, waar $V^\perp$ altijd werkt
(\cref{ch:b3:hilbert}), zijn de uitzondering, niet de regel.
\end{exercise}

\section{Probleem: de dualiteit van de
\texorpdfstring{$\ell^p$}{lp}-ruimten}

\begin{problem}[{Weekendopgave --- $(\ell^p)' = \ell^q$,
reflexiviteit, en het vreemde van $\ell^\infty$}]
\label{pb:b3:banach:1}
Leg $1 < p < \infty$ vast en zij $q$ de toegevoegde exponent, met
$\frac1p + \frac1q = 1$. De paring is overal $\langle x, y\rangle
= \sum_n x_ny_n$.

\medskip
\noindent\textbf{Deel I --- Hölder en Minkowski voor rijen.}
\begin{enumerate}
  \item (Ongelijkheid van Young) Toon voor $a, b \geq 0$ aan dat
        $ab \leq \frac{a^p}p + \frac{b^q}q$, met de concaviteit
        van $\log$ of door $t \mapsto \frac{t^p}p + \frac1q - t$
        te bestuderen.
  \item (Hölder) Leid af dat $\abs{\langle x, y\rangle} \leq
        \norm x_p\norm y_q$ voor $x \in \ell^p$ en $y \in \ell^q$;
        bepaal het gelijkheidsgeval.
  \item (Minkowski) Leid de driehoeksongelijkheid voor
        $\norm\cdot_p$ af. \emph{(Schrijf $\abs{x_n + y_n}^p \leq
        \abs{x_n}\,\abs{x_n{+}y_n}^{p-1} +
        \abs{y_n}\,\abs{x_n{+}y_n}^{p-1}$ en pas Hölder toe op
        elke term.)}
  \item Bewijs dat $\ell^p$ \omterm{def:b3:complete:complete}{volledig} is en dat de eindige rijen er
        \omterm{def:b3:topology:interior}{dicht} in liggen.
\end{enumerate}

\medskip
\noindent\textbf{Deel II --- De dualiteit $(\ell^p)' =
\ell^q$.}
\begin{enumerate}[resume]
  \item Toon voor $y \in \ell^q$ aan dat $\Lambda_y(x) = \langle
        x, y\rangle$ een $\Lambda_y \in (\ell^p)'$ definieert met
        $\norm{\Lambda_y} \leq \norm y_q$, en, door te toetsen op
        $x_n = \abs{y_n}^{q-1}\operatorname{sign}(y_n)$ (passend
        afgeknot en genormaliseerd), dat $\norm{\Lambda_y} = \norm
        y_q$.
  \item Stel omgekeerd $y_n = \Lambda(e_n)$ voor een gegeven
        $\Lambda \in (\ell^p)'$; toon aan dat $y \in \ell^q$ met
        $\norm y_q \leq \norm\Lambda$ \emph{(toets op afknottingen
        zoals in vraag 5 en laat de afknotlengte groeien)}, en
        besluit dat $\Lambda = \Lambda_y$: de afbeelding $y
        \mapsto \Lambda_y$ is een isometrisch isomorfisme $\ell^q
        \to (\ell^p)'$.
  \item Leid af dat $\ell^p$ \emph{\omterm{rem:b3:banach:bidual}{reflexief}} is voor $1 < p <
        \infty$: door de twee dualiteiten samen te stellen komt
        elk element van $(\ell^p)''$ uit $\ell^p$; ga zorgvuldig
        na dat de samenstelling de kanonieke $J$ is.
\end{enumerate}

\medskip
\noindent\textbf{Deel III --- $\ell^1$ en $\ell^\infty$ zijn
andere beesten.}
\begin{enumerate}[resume]
  \item Toon aan dat $\ell^p$ ($1 \leq p < \infty$) en $c_0$
        separabel zijn, maar $\ell^\infty$ niet. \emph{(De
        overaftelbaar vele indicatorrijen van deelverzamelingen
        van $\N$ liggen paarsgewijs op afstand $1$.)}
  \item Leid uit \cref{exo:b3:banach:10} af dat $(\ell^1)' \cong
        \ell^\infty$ maar $(\ell^\infty)' \not\cong \ell^1$:
        $\ell^1$ is \emph{niet} \omterm{rem:b3:banach:bidual}{reflexief}. \emph{(Was
        $(\ell^\infty)'$ gelijk aan $\ell^1$, dan was ze separabel
        en dus $\ell^\infty$ ook.)}
  \item (Een banachlimiet, expliciet) Stel op $\ell^\infty_\R$:
        $p(x) = \limsup_n \frac{x_1 + \dots + x_n}{n}$. Toon aan
        dat $p$ sublineair is, en dat op de deelruimte $c$ van de
        convergente rijen $\mathrm{LIM}(x) = \lim x$ voldoet aan
        $\mathrm{LIM} \leq p$. Zet met Hahn--Banach voort tot
        $\mathrm{LIM} \colon \ell^\infty_\R \to \R$ en toon aan:
        $\mathrm{LIM}$ is positief (uit $x \geq 0$ volgt
        $\mathrm{LIM}(x) \geq 0$), verschuivingsinvariant
        ($\mathrm{LIM}(x_2, x_3, \dots) = \mathrm{LIM}(x)$), zet
        de limiet voort, en voldoet aan $\liminf x \leq
        \mathrm{LIM}(x)\leq \limsup x$.
  \item Toon aan dat zo'n $\mathrm{LIM}$, opgevat in
        $(\ell^\infty)'$, \emph{niet} van de vorm $\Lambda_y$ is
        voor enige $y \in \ell^1$; besluit opnieuw dat
        $(\ell^\infty)' \neq \ell^1$. \emph{(Evalueer op de
        eenheidsrijen $e_n$, en daarna op de constante rij $1$.)}
  \item Evalueer $\mathrm{LIM}$ op $(0,1,0,1,\dots)$, en toon aan
        dat er geen verschuivingsinvariante
        \emph{multiplicatieve} voortzetting van de limiet kan
        bestaan \emph{(beschouw $x = (0,1,0,1,\dots)$ en $x\cdot
        Sx$, waarbij $S$ de verschuiving is)}.
\end{enumerate}

\medskip
\noindent\textbf{Deel IV --- Slot: waarom reflexiviteit ertoe
doet.}
\begin{enumerate}[resume]
  \item Toon met \cref{cor:b3:banach:bslimits} en vraag 6 aan dat
        elke begrensde rij van $\ell^p$ ($1 < p < \infty$) een
        deelrij $(x^{(k)})$ heeft die \emph{zwak} convergeert:
        $\Lambda(x^{(k)})$ convergeert voor elke $\Lambda \in
        (\ell^p)'$. \emph{(Diagonaalextractie over de aftelbaar
        veel coördinaten; bepaal de zwakke limiet in $\ell^p$ met
        de uniforme begrensdheid van de normen en Hölder.)} Toon
        aan de hand van een voorbeeld ($e_n$ in $\ell^1$, tegen
        goed gekozen elementen van $\ell^\infty$) aan dat dit in
        $\ell^1$ faalt: zwakke \omterm{def:b3:topology:compact}{compactheid} is een voorrecht van
        \omterm{rem:b3:banach:bidual}{reflexieve ruimten}.
\end{enumerate}

\medskip
\noindent\textbf{Deel V --- De zwakke \omterm{def:b3:topology:topology}{topologie} aan het werk, en
de verrassing van Schur.} Schrijf $x^{(k)} \rightharpoonup x$ in
een genormeerde ruimte $E$ (\emph{zwakke convergentie}) wanneer
$\Lambda(x^{(k)}) \to \Lambda(x)$ voor elke $\Lambda \in E'$.
\begin{enumerate}[resume]
  \item Maak de inventaris af: toon isometrisch aan dat $(c_0)'
        \cong \ell^1$, volgens het schema van de vragen 5--6 (wat
        neemt de plaats van de toetsrijen in?). Stel de keten $c_0
        \to \ell^1 \to \ell^\infty \to \dots$ van opeenvolgende
        \omterm{def:b3:banach:operator}{duale ruimten} samen en markeer waar de reflexiviteit
        faalt.
  \item Toon aan dat elke zwak convergente rij van een
        banachruimte begrensd is: bekijk de $x^{(k)}$ via de
        kanonieke inbedding $J$ als functionalen op $E'$ en pas
        Banach--Steinhaus toe
        (\cref{thm:b3:banach:banachsteinhaus}) --- op welke
        banachruimte, en waarom is de \omterm{def:b3:complete:complete}{volledigheid} daar
        beschikbaar?
  \item Toon aan dat in $\ell^p$ met $1 < p < \infty$ geldt:
        $x^{(k)} \rightharpoonup x$ dan en slechts dan als
        $\sup_k\norm{x^{(k)}}_p < \infty$ en $x^{(k)}_n \to x_n$
        voor elke coördinaat $n$ \emph{(één richting gebruikt
        Banach--Steinhaus via de kanonieke inbedding; benader voor
        de andere $y \in \ell^q$ door eindige rijen)}. Leid af dat
        $e_k \rightharpoonup 0$ in $\ell^2$ terwijl
        $\norm{e_k}_2 = 1$: zwakke limieten kunnen massa
        verliezen.
  \item Toon aan dat de norm zwak onderhalfcontinu is: uit
        $x^{(k)} \rightharpoonup x$ volgt $\norm x \leq
        \liminf_k\,\norm{x^{(k)}}$ \emph{(kies een normerende
        functionaal voor $x$, \cref{cor:b3:banach:hbnormed})}.
  \item (Radon--Riesz in $\ell^2$) Toon aan dat in $\ell^2$ zwakke
        convergentie samen met convergentie van de normen
        normconvergentie impliceert \emph{(werk $\norm{x^{(k)} -
        x}_2^2$ uit)}. Geef een tegenvoorbeeld voor dezelfde
        uitspraak zonder de normhypothese.
  \item (Schur, stap 1) Zij $x^{(k)} \rightharpoonup 0$ in
        $\ell^1$ en stel, uit het ongerijmde, dat
        $\norm{x^{(k)}}_1 \geq \delta > 0$ langs een deelrij. Toon
        eerst aan dat $x^{(k)}_n \to 0$ voor elke $n$ (welke
        functionalen?), en construeer daarna recursief indices
        $k_1 < k_2 < \cdots$ en gehele getallen $0 = N_0 < N_1 <
        N_2 < \cdots$ zodanig dat de massa van $x^{(k_j)}$ zich
        concentreert op het blok $B_j =
        \intoc{N_{j-1}}{N_j}$:
        \[
        \sum_{n \in B_j}\bigl|x^{(k_j)}_n\bigr| \geq
        \norm{x^{(k_j)}}_1 - \frac\delta{10} .
        \]
  \item (Schur, stap 2) Definieer $y \in \ell^\infty$ door $y_n =
        \operatorname{sign}\bigl(x^{(k_j)}_n\bigr)$ voor $n \in
        B_j$. Toon aan dat $\langle x^{(k_j)}, y\rangle \geq
        \norm{x^{(k_j)}}_1 - \frac{2\delta}{10} \geq
        \frac{8\delta}{10}$ en leid een tegenspraak af met
        $x^{(k)} \rightharpoonup 0$. Besluit de \emph{stelling van
        Schur}: in $\ell^1$ convergeren zwak convergente rijen in
        norm.
  \item Leid af dat $(e_k)$ in $\ell^1$ geen zwak convergente
        deelrij heeft (haar enige kandidaat-limiet is $0$, per
        coördinaat --- pas dan Schur toe), waarmee het falen van
        de zwakke \omterm{def:b3:topology:compact}{compactheid} uit vraag 13 wordt teruggevonden; en
        los de schijnbare paradox op: in $\ell^1$ vallen zwakke en
        normconvergentie van \emph{rijen} samen, en toch
        verschillen de zwakke en de norm\emph{\omterm{def:b3:topology:topology}{topologie}} en zijn
        begrensde verzamelingen nog steeds niet zwak rijcompact
        --- geen tegenspraak, alleen het falen van de
        reflexiviteit.
  \item (Overzichtstabel) Stel voor $E \in \{c_0,\ \ell^1,\
        \ell^p\ (1{<}p{<}\infty),\ \ell^\infty\}$ een tabel op
        met: de \omterm{def:b3:banach:operator}{duale ruimte}; separabiliteit; reflexiviteit; of
        begrensde rijen zwak convergente deelrijen toelaten; en
        één kenmerkende eigenschap van elke ruimte, in één regel
        verantwoord vanuit dit probleem.
\end{enumerate}

\medskip
\noindent\textbf{Deel VI --- Aanvullingen: dichtstbijzijnde
punten, gemiddelde convergentie, de waarde van een banachlimiet.}
\begin{enumerate}[resume]
  \item (Dichtstbijzijnde punten: een dividend van de
        reflexiviteit) Zij $F$ een gesloten deelruimte van
        $\ell^p$ ($1 < p < \infty$) en $x \in \ell^p$. Toon aan
        dat $d = \operatorname{dist}(x, F)$ \emph{bereikt} wordt:
        haal uit een minimaliserende rij een zwak convergente
        deelrij (vraag 13), houd de zwakke limiet binnen $F$ door
        met Hahn--Banach een functionaal te bouwen die op $F$
        verdwijnt maar niet in een punt daarbuiten, en besluit met
        vraag 17. Toon daarna aan dat dit voorrecht niet
        universeel is: bewijs voor $\Lambda(x) = \sum_n2^{-n}x_n$
        op $c_0$ dat $\norm\Lambda = 1$ niet op de eenheidsbal
        bereikt wordt, stel de afstandsformule
        $\operatorname{dist}(x, \ker\Lambda) = \abs{\Lambda(x)}$
        op, en leid af dat geen enkele $x \notin \ker\Lambda$ een
        dichtstbijzijnd punt in het gesloten hypervlak
        $\ker\Lambda$ heeft.
  \item (Banach--Saks in $\ell^2$) Zij $x^{(k)} \rightharpoonup 0$
        in $\ell^2_{\R}$ met $\norm{x^{(k)}}_2 \leq C$. Construeer
        een deelrij $(y_j)$ met $\abs{\langle y_i, y_j\rangle}
        \leq \frac1j$ voor alle $i < j$, en leid af dat
        \[
        \Bigl\lVert\frac{y_1 + \dots + y_m}m\Bigr\rVert_2^2
        \leq \frac{C^2 + 2}m \longrightarrow 0 :
        \]
        na extractie convergeren de cesàromiddelen dus in
        \emph{norm}. Controleer dat op $(e_k)$, waarvan de
        middelen norm $\frac1{\sqrt m}$ hebben: zwakke
        convergentie, nutteloos voor de rij zelf (vraag 16), wordt
        normconvergentie voor de gemiddelden.
  \item (De waarde van een banachlimiet) Zij $L$ een willekeurige
        banachlimiet (vraag 10) en $A_mx = \frac1m(x + Sx + \dots
        + S^{m-1}x)$. Toon aan dat $L(A_mx) = L(x)$ en $\liminf
        A_mx \leq L(x) \leq \limsup A_mx$; leid af dat \emph{alle}
        banachlimieten op periodieke rijen samenvallen, met als
        waarde het gemiddelde over een periode --- $\frac13$ op
        $(1, 0, 0, 1, 0, 0, \dots)$, in overeenstemming met de
        $\frac12$ uit vraag 12. Toon daarna aan dat die
        overeenstemming in het algemeen faalt: toon voor de
        blokrij $x$ die gelijk is aan $1$ op
        $\intoc{3^{j-1}}{3^j}$ voor even $j$ en $0$ elders, aan
        dat de cesàromiddelen tussen $\leq \frac13$ en $\geq
        \frac23$ schommelen, en bouw twee banachlimieten $L_\pm$
        met $L_-(x) \leq \frac13 < \frac23 \leq L_+(x)$
        \emph{(zet voort vanaf $c \oplus \R x$ met de uiterste
        toegelaten waarden $\pm$: ga na dat $\Lambda(y + tx) =
        \lim y + t\,p(x)$ door de sublineaire $p$ van vraag 10
        wordt gedomineerd)}.
\end{enumerate}
\end{problem}
